% !TEX root = ../main.tex
\section{Exponential Growth}

\begin{note-block}
The data $x[n] = r^n + w[n]$ for $n = 0, 1, \ldots , N - 1$ are observed, where $w[n]$ is WGN with known variance $\sigma^2$ and $r > 0$ is to be estimated.
\begin{enumerate}[label={(\alph*)}]
    \item Find the CRLB for $r$.
    \item Does an efficient estimator exist and if so find its variance?
\end{enumerate}
\end{note-block}

We arrive at the Log-Likelihood function in the same way as for \autoref{eq:log-likelihood-A}, but for the parameter $r$
\begin{align}
    \ln(p(\vec{x};r)) = -\frac{N}{2}\ln(2\pi \sigma^2_w) -\frac{1}{2\sigma_w^2}\sum_{n=0}^{N-1} (x[n] - r^n)^2 \label{eq:log-likelihood-r}
\end{align}
First Derivative
\begin{align}
    \pd[\ln(p(\vec{x};r))]{r}
    &= -\frac{1}{2\sigma_w^2}\sum_{n=0}^{N-1} 2(x[n] - r^n) (-nr^{n-1})\notag \\
    &= \frac{1}{\sigma_w^2}\sum_{n=0}^{N-1} (nx[n]r^{n-1} - nr^{2n-1}) \label{eq:dL-r}
\end{align}
Second Derivative of the Log Likelihood Function
\begin{align}
    \pnd[\ln(p(\vec{x};r))]{2}{r}
    &= \frac{1}{\sigma_w^2} \sum_{n=0}^{N-1} (n(n-1)x[n]r^{n-2} - n(2n-1)r^{2n-2})
\end{align}
Fisher Information
\begin{align*}
    I(r) = -\E{\pnd[\ln(p(\vec{x};r))]{2}{r}}
    = -\frac{1}{\sigma_w^2}
    \E{\sum_{n=0}^{N-1} (n(n-1)x[n]r^{n-2} - n(2n-1)r^{2n-2})}
\end{align*}
All terms, except $x[n]$ are deterministic. From the linearity of the expected value operator and $\E{x[n]}=r^n$, because $w[n]$ has zero mean, follows.
\begin{align}
    I(r) 
    &= \frac{1}{\sigma_w^2} \sum_{n=0}^{N-1} (n(2n-1)r^{2n-2} - n(n-1)\E{x[n]}r^{n-2}) \notag \\
    &= \frac{1}{\sigma_w^2} \sum_{n=0}^{N-1} (n(2n-1)r^{2n-2} - n(n-1)r^{2n-2}) \notag \\
    &= \frac{1}{\sigma_w^2} \sum_{n=0}^{N-1} (n(2n-1-n+1)r^{2n-2}) \notag \\
    &= \frac{1}{\sigma_w^2} \sum_{n=0}^{N-1} n^2r^{2n-2} \label{eq:fisher-information-r}
\end{align}

The \acs{CRLB} is
\begin{boxed-block}
\begin{equation}
\Var{\hat r} \geq \frac{1}{I(r)} = \frac{\sigma_w^2}{\sum_{n=0}^{N-1} n^2r^{2n-2}}
\end{equation}
\end{boxed-block}

\subsection{Efficient Estimator}

Again, we want to find an estimator $g(\vec{x})$ such that
\begin{equation}
    \pd[\ln(p(\vec{x};r))]{r}=I(r)(g(\vec{x})-r)
\end{equation}

Using \autoref{eq:dL-r} on the left side and the Fisher Information from \autoref{eq:fisher-information-r}, we get
\begin{align*}
    \frac{1}{\sigma_w^2}\sum_{n=0}^{N-1} (nx[n]r^{n-1} - nr^{2n-1})
    &\overset{!}{=} \frac{1}{\sigma_w^2} \left(\sum_{n=0}^{N-1} n^2r^{2n-2}\right)(g(\vec{x}) - r) \\
\end{align*}
Solving for $g(\vec{x})$ would give
\[
g_{N=1}(\vec x) = \frac{\sum_{n=0}^{N-1} (nx[n]r^{n-1} - nr^{2n-1})}{\sum_{n=0}^{N-1} (nr^{2n-1})^2} + r
\]
which is a function of the true, unknown value of $r$, so the estimator doesn't work. An exception for this, is when $N = 2$:
\[
g(\vec{x}) = \frac{x[1]r^0 - r}{(r^0)^2} - r = x[1]
\]
So for this case, $x[1]=\hat r$ is the best estimate for $r$. The variance for this estimator is simply $\sigma_w^2$.
